% !TeX root = ../PIFB2.tex
\section{What Formation Theory Establishes}\label{sec:participatory}

The companion theory establishes pairing selection in a declared four-agent model. The result distinguishes learned from frozen source priors. Its scope must be separated from physical object formation and an indefinitely recursive hierarchy.

\subsection{The fixed-temperature selection result}
The model has a homogeneous cycle of four agents at $c_*$. Its configuration consists of four scalar belief coordinates, four model coordinates and one shared scalar translation on each unoriented edge. With a chosen orientation of edge $e=i\to i+1$, the comparison cost is
\begin{equation}
D_e=\frac12\left[(k_{i+1}-k_i-w_e)^2+(m_{i+1}-m_i-w_e)^2\right].
\label{eq:formation_cycle_cost}
\end{equation}
Both directions of source selection use the same edge variable and cost. A Gaussian reference confines the configuration, a stochastic dynamics evolves it at temperature $T>0$, and source assignments are refreshed using their conditional law. The product source-prior rows start uniform and learn from the actual recognition law through
\begin{equation}
\dot\pi_{ij}=\epsilon_L\bigl(q_t(J_i=j)-\pi_{ij}\bigr).
\label{eq:formation_source_learning}
\end{equation}
Write $S$ for the occupied unoriented edge set induced by the source assignments $J$. These ingredients, including the reference and kinetics, are part of the model. The statement is not transferable to arbitrary links or arbitrary source updates.

For the initial law $\mathcal N(m_0,I_{12})$ times uniform assignments, the companion theorem proves: there is $T_0>0$ such that for every fixed $T\in(0,T_0]$ a finite constant $C_T$ exists and, for Lebesgue-almost every fixed $m_0$, one perfect matching $M$ satisfies
\begin{equation}
\lim_{\epsilon_L\to0}
\sup_{t_L\leq t\leq2t_L}q_t^{\epsilon_L}(S\ne M)=0,
\qquad
t_L=C_T\epsilon_L^{-1}\log(1/\epsilon_L).
\label{eq:formation_selection_window}
\end{equation}
At frozen uniform rows, for every $T>0$ the long-time difference of the two matching probabilities instead tends to zero~\cite[Ch.~15 and the selection proof]{DennisGaugeVFE2026}.

The proof combines frozen-law entropy contraction, a nonvanishing matching-mode imprint from almost every initial mean, the small-temperature landscape and controlled escape on the logarithmic slow horizon. The initial mean supplies the asymmetry amplified by learning. The statement neither starts from an exactly symmetry-invariant law and spontaneously chooses a deterministic sign, nor takes a joint temperature--mobility limit. It concerns the actual law uniformly over the stated window, not the probability that a single trajectory shows a chosen pair once.

\subsection{Selection, inheritance and the next episode}
A separate seeded-formation result can show an episode with positive probability even when source learning is idle. Selection against a symmetric frozen control therefore tests a stronger mechanism than observing a coherent cluster. The theorem closes that mechanism for its declared family.

Three questions remain distinct. A selected pairing must first admit a retained parent state and outgoing kernels. Its law must then persist or carry the necessary conditional memory during a second episode. Finally, the resulting parent dynamics must itself support another formation step without replacing inherited parameters by newly chosen ones. The theory has conditional parent constructions and supplied finite-depth examples, but those do not settle the general recursive problem. In particular, a slow model record can account for persistence while failing to cause the next coalescence.

Rank-deficient selective communication changes this problem. Public agreement can hide private variables that later affect parent interactions. The descent square~\eqref{eq:parent_link_square}, generator condition~\eqref{eq:generator_descent} and conditional-memory state identify what has to be retained. An extension that merely adds low-rank notation to the four-cycle theorem would not prove formation with learned selective interfaces.

\subsection{A physical comparison requires an external target}
The theorem is a mathematical result about the model. A physical claim would additionally identify a target system, observable records, an inferred state and channel family, parameter uncertainty and a prediction not fitted at the coarse scale. Appropriate contrasts include frozen source learning, altered interface rank, and parent predictions with and without retained memory, provided each comparison preserves the stated model assumptions.

No arbitrary detector threshold, seed count or pass percentage is elevated to a physical law here. A validation design must derive its resolution and uncertainty from the target measurement and the particular claim. Failure to match that target would reject the physical identification or approximation, not the conditional algebra of Bayesian inference. Conversely, reproducing a declared stochastic construction would not establish that nature implements it.

